\section{Invariant expansion}

\label{app:G}

For a traceless symmetric 3 by 3 matrix, Cayley--Hamilton gives $h^3-I_2h/2-I_3I/3=0$. Multiplication by $h$ and tracing yields $\mathrm{tr}(h^4)=I_2^2/2$. Expanding the exponentials,

\begin{equation}\tag{G1}\label{eq:G1}
\begin{aligned}
\mathrm{tr}(e^{2h})&=3+2I_2+\frac43I_3+\frac13I_2^2+O(h^5),\\
\mathrm{tr}(e^{4h})&=3+8I_2+\frac{32}{3}I_3+\frac{16}{3}I_2^2+O(h^5).
\end{aligned}
\end{equation}

Taking twice the second expression minus the square of the first gives (I1). Since $h$ is symmetric, $|e^h\omega|^2=\omega^{\mathsf T}e^{2h}\omega$, proving the twist--tensor terms as well. The scalar coefficients in these expansions depend on $r,L,\eta,V$ through (M3); expanding those coefficients adds higher-order scalar couplings.

Rotation invariance would allow $h_{ij}u_iv_j$ for two independent vectors $u,v$. Whether such an invariant occurs is decided by the actual action and the meaning of its variables. Here the mixing coordinates parametrize a coframe rotation and drop out of the scalar action. Their absence from (I1) is consequently an action-specific result. The derivative invariant $Q$ in Section 4 belongs to the separately specified circle-dependent vector family.
